% FILE: 01_research_question.tex
% Project: otel-weaver
% Role: telemetry-to-process-evidence-ingestion-layer

\section{Research Question}
\label{sec:otel-weaver-rq}

\subsection{Problem Statement}
\label{sec:otel-weaver-rq-problem}

Contemporary distributed systems instrumentation via OpenTelemetry produces a high-volume stream of
spans, metrics, and log records. Practitioners and platform teams routinely treat this stream as
process evidence: dashboards display ``process health,'' schema validation passes are cited as
compliance proofs, and alert thresholds on span latency are described as SLA enforcement. The
problem \texttt{otel-weaver} investigates is that none of these treatments is formally defensible.
A span is a measurement artifact. A schema validation result is a structural check. Neither
constitutes a process conformance finding because neither references a declared process model, a
lawful lifecycle state, or a cryptographically bound receipt chain.

The failure mode this creates is not theoretical. When observability signals are silently promoted
to process evidence---without a typed admission gate, without a loss report, and without a witness
seal---the downstream process-intelligence layer receives feedstock it cannot distinguish from
ground-truth evidence. Conformance checking performed against such input produces fitness scores
that measure instrument coverage rather than process compliance. The resulting verdicts are
category errors, not process findings.

The research problem is therefore: \emph{What formal boundary is required between OpenTelemetry
observability data and process-intelligence evidence, and what type-system mechanisms enforce that
boundary without requiring the ingestion layer to perform conformance checking itself?}

\subsection{Old-Regime Assumption Challenged}
\label{sec:otel-weaver-rq-old-regime}

The old-regime assumption \texttt{otel-weaver} directly challenges is encoded in doctrine constant
\texttt{NO\_DASHBOARD\_TRUTH\_001}:

\begin{quote}
  Observability dashboards display instrument readings. They do not display process truth.
  No dashboard output, alert, or visualization constitutes a process conformance verdict.
\end{quote}

This assumption is pervasive in the industry. Platform engineering teams speak of ``observing''
processes, ``monitoring'' compliance, and ``detecting'' process drift through metric anomalies.
The doctrine file \texttt{no-dashboard-truth.md} names this assumption explicitly and refuses it:
a dashboard is a visualization of feedstock, not a process oracle.

The companion assumption challenged by doctrine constant
\texttt{OBSERVABILITY\_VS\_PROCESS\_INTELLIGENCE\_001} is that observability-by-design and
process intelligence are architecturally equivalent or substitutable. The doctrine establishes
that they are categorically distinct: observability captures what the system \emph{emitted};
process intelligence certifies what the process \emph{did}. The former is necessary but not
sufficient for the latter.

A third assumption, challenged by \texttt{WEAVER\_FINDING\_NOT\_RECEIPT\_001}, is that a Weaver
\texttt{registry live-check} finding---a real-time feedstock validation result---constitutes a
receipt in the foundry sense. It does not. A live-check finding is an instrument reading about
schema conformance at ingestion time. A receipt is a BLAKE3-chained signed assertion about
process-lifecycle compliance at verdict time. These are not the same object and must not be
conflated in any downstream claim.

\subsection{Hypothesis}
\label{sec:otel-weaver-rq-hypothesis}

The central hypothesis of the \texttt{otel-weaver} project is:

\begin{quote}
  A typed ingestion gate enforcing \texttt{Admission<T,\,W>} / \texttt{Refusal} at the
  OTel-to-process-evidence boundary is both necessary and sufficient to prevent category collapse,
  provided the gate is backed by: (a) a formal doctrine layer naming what telemetry is not,
  (b) a loss-report requirement making projection decisions explicit, and (c) a witness-seal
  requirement binding each admitted payload to a lattice-bounded \texttt{ActivityWitness}.
\end{quote}

The hypothesis is testable: if any span can reach a downstream process-intelligence consumer
without traversing the gate, the gate is insufficient. If any gate refusal fails to emit a
structured \texttt{Refusal} with a BLAKE3 digest, the gate is unsound. Experiment \texttt{exp-003}
(\textit{live-check-to-refusal}) provides the primary test surface for this hypothesis, with
inline test cases \texttt{test\_successful\_admission} and \texttt{test\_refusal\_missing\_id}
specified in its README.

\subsection{Success Criteria}
\label{sec:otel-weaver-rq-success}

Success for the \texttt{otel-weaver} research component is defined by three ordered criteria,
corresponding to the three \texttt{ggen}-signed checkpoints:

\begin{enumerate}
  \item \textbf{Design gate} (\texttt{PARTIAL\_001}, SHA-256: \texttt{d4b844c4\ldots}):
    All six doctrine files authored, all five mapping tables specified as ACTIVE, and
    \texttt{process.pi} registry defined. This checkpoint certifies conceptual completeness.
  \item \textbf{Full asset completion} (\texttt{ALIVE\_001}, SHA-256: \texttt{e6bc59b8\ldots}):
    All five experiment READMEs complete with inline Rust implementation specifications,
    \texttt{ggen} manifests manufacturing the BLAKE3 receipt chain, and the Ingestion Covenant
    (Admit-Don't-Assume, Loss Reports, Witness Seal) formally declared.
  \item \textbf{Runtime integration gate} (\texttt{RUNTIME\_001}, SHA-256: \texttt{86ef45b3\ldots}):
    At least one passing test in \texttt{weaver\_integration\_tests.rs} demonstrating that the
    ingestion pipeline connects to \texttt{wasm4pm} and \texttt{wasm4pm-compat} without
    category collapse. The checkpoint narrative cites 1 passed / 0 failed in this file and
    62 total passing tests across the dependent packages.
\end{enumerate}

The project's current status is PARTIAL: all three checkpoints have been signed by \texttt{ggen},
but the test execution evidence for \texttt{RUNTIME\_001} exists as narrated text within the
checkpoint file rather than as materialized compiled artifacts or test runner output on disk.
Full ALIVE certification requires materialized runtime evidence.
